\section{Hermitian 矩阵与复谱定理}
\label{sec:03-Linear-Algebra/08-Complex-Matrices-and-Fourier-Structure/02}

一组振动传感器同时采样，你把每个通道在某个频率上的 Fourier 系数收成复向量 \(z^{(r)}\)，\(r\) 标记时间窗，然后照统计里的老办法做二阶矩。接下来要回答的问题很实际：按权重 \(u\) 把各通道线性叠加之后，这个组合信号在该频率上的能量有多大？

如果按实数里的习惯写成 \(C=N^{-1}\sum_r z^{(r)}(z^{(r)})^{\trans}\)，再算 \(u^{\trans}Cu\)，得到的是一个普通复数——可能带虚部，可能实部为负。能量是可测量的物理量，不该长成这样。换成共轭转置，

\[
C=\frac1N\sum_{r=1}^N z^{(r)}(z^{(r)})^{\conj},
\qquad
u^{\conj}Cu=\frac1N\sum_{r=1}^N\abs{(z^{(r)})^{\conj}u}^2\ge0,
\]

数值自动落回非负实轴。差别只在一个共轭符号上，可它决定了这个矩阵能不能当协方差用。这样得到的 \(C\) 满足 \(C=C^{\conj}\)，对角元记录各通道功率，非对角元同时保存共同振动的强度与相位。

\begin{center}
\begin{tikzpicture}[x=1cm,y=1cm,font=\footnotesize,>=stealth]
  \draw[->,black!45] (-0.6,0) -- (4.3,0) node[below right] {\(\operatorname{Re}\)};
  \draw[->,black!45] (0,-1.5) -- (0,1.6) node[left] {\(\operatorname{Im}\)};
  \foreach \p in {0.35,0.72,1.15,1.6,2.05,2.4,2.85,3.3,3.75} {
    \fill[blue!65!black] (\p,0) circle (2pt);
  }
  \node at (1.9,-2.1) {\(H=H^{\conj}\)};
  \begin{scope}[shift={(6.6,0)}]
    \draw[->,black!45] (-1.9,0) -- (2.3,0) node[below right] {\(\operatorname{Re}\)};
    \draw[->,black!45] (0,-1.5) -- (0,1.6) node[left] {\(\operatorname{Im}\)};
    \foreach \p/\q in {0.9/0.7, -0.5/1.1, 1.6/-0.4, -1.2/-0.8, 0.2/-1.2, 1.1/1.25, -1.5/0.35, 0.55/-0.65, 1.85/0.95} {
      \fill[red!60!black] (\p,\q) circle (2pt);
    }
    \node at (0.2,-2.1) {\(S=S^{\trans}\)，但 \(S\ne S^{\conj}\)};
  \end{scope}
\end{tikzpicture}

\emph{同一批单位向量 \(x\) 代进二次型 \(x^{\conj}Ax\)，画出结果在复平面上的位置。左边的 Hermitian 条件把所有取值钉在实轴上，能量、方差、测量值这些概念才有意义；右边只是``沿对角线镜像相等''的复对称矩阵，取值散落一片，最大最小之类的话根本无从说起。}
\end{center}

\subsection{\texorpdfstring{把 \(A^{\trans}=A\) 改写成 \(A^{\conj}=A\)}{把 A 转置等于 A 改写成 A 共轭转置等于 A}}

实对称矩阵的特殊地位不来自视觉上的对角线对称，而来自它等于自己关于内积的伴随。上一节已经把伴随从 \(A^{\trans}\) 换成了 \(A^{\conj}\)，所以对应条件是

\[
H=H^{\conj},
\]

称为 Hermitian（厄米）矩阵。逐项写就是 \(h_{ij}=\overline{h_{ji}}\)，于是对角元必为实数，虚部只能寄生在成对的非对角位置上。

顺手排除一个替代品：\(H=H^{\trans}\) 这个不带共轭的条件不能用。矩阵

\[
S=\begin{bmatrix}1&i\\i&-1\end{bmatrix}
\]

满足 \(S=S^{\trans}\)，看着很对称，可 \(S^2=0\)。它的特征值全是 \(0\)，独立特征向量只有一个，连对角化都做不到，更谈不上正交特征基。镜像对称是坐标层面的巧合，共轭对称才对应内积层面的结构。

\subsection{共轭在证明的哪两步救了场}

先看二次型。对任意 \(x\)，标量 \(x^{\conj}Hx\) 的共轭等于它自己的共轭转置，于是

\[
\overline{x^{\conj}Hx}=(x^{\conj}Hx)^{\conj}=x^{\conj}H^{\conj}x=x^{\conj}Hx,
\]

一个数等于自己的共轭，只能是实数。这里共轭出现了两次：一次是把标量共轭改写成共轭转置，一次是用 \(H^{\conj}=H\) 消掉中间的矩阵。少任何一次，链条就断在中途。

接着是特征值。设 \(Hv=\lambda v\)、\(v\ne0\)，两边左乘 \(v^{\conj}\)：

\[
\lambda=\frac{v^{\conj}Hv}{v^{\conj}v}.
\]

分子刚证明是实数，分母 \(v^{\conj}v=\norm{v}^2>0\) 也是实数——这一步靠的正是上一节修好的正定性，用 \(v^{\trans}v\) 就可能得到零分母。所以 \(\lambda\) 实。

正交性同理。若 \(Hv=\lambda v\)、\(Hw=\mu w\) 且 \(\lambda\ne\mu\)，由 \(\lambda\) 为实数有 \((\lambda v)^{\conj}=\lambda v^{\conj}\)，代入伴随关系：

\[
\lambda\, v^{\conj}w=(Hv)^{\conj}w=v^{\conj}Hw=\mu\, v^{\conj}w,
\]

两边相减得 \((\lambda-\mu)v^{\conj}w=0\)，故 \(v^{\conj}w=0\)。注意这里也用了``\(\lambda\) 是实数''，否则左边会冒出一个 \(\overline\lambda\)，等式两侧不再可比。

\begin{theorem}[复谱定理]
方阵 \(H\) 满足 \(H=H^{\conj}\)，当且仅当存在酉矩阵 \(U\) 与实对角矩阵 \(\Lambda\) 使 \(H=U\Lambda U^{\conj}\)。
\end{theorem}

一个方向一行就够：对 \(U\Lambda U^{\conj}\) 取共轭转置，由 \(\Lambda^{\conj}=\Lambda\) 立刻回到自身。另一个方向的骨架与实情形完全相同，对维数归纳：代数封闭保证至少存在一个特征值 \(\lambda_1\) 与单位特征向量 \(u_1\)；由伴随关系可验证 \(u_1\) 的正交补是 \(H\) 的不变子空间，把 \(H\) 限制上去仍是 Hermitian，维数降一，归纳假设接手。实情形里这个骨架有一处要额外操心——特征值会不会跑到复数去，从而特征向量不在实空间里；换到复空间反而不必操心，因为标量域已经封闭。

于是实对称矩阵的全部结论原样过户：一组标准正交的特征基、二次型在特征坐标下拆成独立平方项、正定等价于特征值全正。变化只有两处，转置写成共轭转置，正交矩阵改叫酉矩阵。

\subsection{一个满是虚数却拥有实谱的例子}

取

\[
H=\begin{bmatrix}2&i\\-i&2\end{bmatrix}.
\]

上下三角的两个元素互为共轭，所以 \(H=H^{\conj}\)。特征多项式 \((2-\lambda)^2-1\) 给出 \(\lambda=1,3\)，都是实数，尽管矩阵里到处是 \(i\)。单位特征向量可取

\[
u_1=\frac1{\sqrt2}\begin{bmatrix}-i\\1\end{bmatrix},
\qquad
u_2=\frac1{\sqrt2}\begin{bmatrix}i\\1\end{bmatrix},
\]

它们满足 \(u_1^{\conj}u_2=0\)。把两列拼成 \(U\)，便有 \(H=U\diag(1,3)U^{\conj}\)。换到这组坐标 \(x=Uy\)，二次型化为 \(x^{\conj}Hx=\abs{y_1}^2+3\abs{y_2}^2\)，非零向量处恒正，可见这个 \(H\) 还是正定的。判定正定不必再看主子式，看特征值就够——这一条也是从实情形整包搬过来的。

\subsection{正规矩阵划出酉对角化的确切边界}

Hermitian 不是能被酉对角化的唯一一类。真正的充要条件更宽：若

\[
A^{\conj}A=AA^{\conj},
\]

称 \(A\) 为正规矩阵（normal matrix），而``\(A\) 可被酉对角化''与``\(A\) 正规''互为等价。Hermitian 矩阵显然正规，但正规的家族要大得多，各类成员的差别恰好体现在特征值落在复平面的哪个位置上。

\begin{center}
\begin{tikzpicture}[x=1.15cm,y=1.15cm,font=\footnotesize,>=stealth]
  \foreach \p/\q in {1.55/1.3, -1.7/0.75, 0.85/-1.55, -1.15/-1.2, 1.9/-0.6, -0.6/1.75} {
    \fill[black!35] (\p,\q) circle (1.8pt);
  }
  \draw[blue!60!black,thick] (0,0) circle (1);
  \draw[red!65!black,very thick] (-2.1,0) -- (2.1,0);
  \draw[green!45!black,very thick] (0,-2.1) -- (0,2.1);
  \draw[->,black!45] (2.1,0) -- (2.5,0) node[below right] {\(\operatorname{Re}\)};
  \draw[->,black!45] (0,2.1) -- (0,2.5) node[left] {\(\operatorname{Im}\)};
  \node[red!65!black,right] at (1.35,-0.32) {Hermitian};
  \node[green!45!black,right] at (0.12,2.1) {skew-Hermitian};
  \node[blue!60!black,below right] at (0.62,-0.72) {酉};
  \node[black!55] at (1.72,1.72) {一般正规};
\end{tikzpicture}

\emph{同一张复平面，四种对称条件把谱分别压在不同的几何对象上：\(H^{\conj}=H\) 压到实轴，\(A^{\conj}=-A\) 压到虚轴，\(U^{\conj}U=I\) 压到单位圆，一般正规矩阵则可以散落任意位置。四类矩阵的共同点是都拥有一组标准正交特征基；不同点只是特征值的取值范围。}
\end{center}

正规不等于 Hermitian。\(\diag(i,2)\) 是正规矩阵，可以酉对角化，但它的谱不在实轴上，二次型 \(x^{\conj}Ax\) 也不总是实数。把这几类摆在一起看，条件与结论的分工就清楚了：正规性负责``存在正交特征基''，Hermitian 性额外负责``特征值是实数、二次型可比较大小''。需要排序、需要谈正定、需要把二次型当能量，就必须要到后面这一条。

越过这条边界之后并非一片空白。任意复方阵都有 Schur 分解 \(A=QTQ^{\conj}\)，\(Q\) 酉、\(T\) 上三角，特征值排在 \(T\) 的对角线上——这是数值特征值算法实际计算的东西。只有正规矩阵能让 \(T\) 进一步退化成对角矩阵。所以``特征多项式在复数中分裂''与``存在标准正交特征基''之间还隔着正规性这一道，代数封闭并不自动兑现几何上的好处。

\subsection{为什么下一步会走到 Fourier}

现在换个方向提问：给定一族彼此可交换的正规矩阵，它们能不能共享同一组特征基？答案是能，而且这正是接下来要用的杠杆。把序列的下标按周期看待，循环移位算子 \(P\) 是酉矩阵，因而正规；由 \(P\) 的各次幂线性组合出来的循环卷积矩阵彼此可交换，也都正规。于是存在一组标准正交向量，把所有这些算子同时对角化。

换句话说，Fourier 变换不是从一堆正弦余弦公式里试出来的。它是一道线性代数问题的解：什么坐标能把全体平移不变算子同时对角化？下一节直接从循环移位的特征方程出发，把这组坐标算出来。
